\chapter{eScript (Electronic Prescription)}
\index{eScript (Electronic Prescription)}

\section*{Definition and Overview}
An electronic prescription, commonly called an eScript, is a digital version of a prescription that a doctor sends to a pharmacy or stores on the person's phone, instead of handing over a paper script. In many countries, eScripts were introduced widely during the COVID-19 pandemic and are now a standard part of health care. The prescription is stored as a digital record with a unique token, such as a QR code or barcode, which the pharmacist scans to dispense the medicine. eScripts are convenient, secure, and reduce the risk of lost or illegible prescriptions.

\section*{Epidemiology}
The use of electronic prescriptions has grown rapidly in many countries. Availability, identity checks, pharmacy systems, and legal requirements vary by jurisdiction. Telehealth has helped drive this growth because a prescription can be sent securely after a remote consultation, but an electronic prescription is still subject to the same clinical and dispensing safeguards as a paper prescription.

\section*{Aetiology and Pathophysiology}
An eScript works through secure digital health systems.
\begin{itemize}
    \item \textbf{Creation:} the prescriber writes the prescription in their clinical software, which generates a digital prescription
    \item \textbf{Delivery:} the prescription can be sent directly to a chosen pharmacy or provided to the patient as a token
    \item \textbf{The token:} a QR code or barcode, or a link, that the patient can store on their phone, in a wallet, or on paper
    \item \textbf{Dispensing:} the pharmacist scans or enters the token, which retrieves the prescription
    \item \textbf{Security:} the system is encrypted and linked to national digital health standards, and tokens are unique to each prescription
    \item \textbf{Repeat prescriptions:} repeats are stored and managed digitally
    \item \textbf{Privacy:} the system is designed so that only the patient and their pharmacy can access the prescription
\end{itemize}

\section*{Clinical Features}
People will encounter eScripts in everyday care.
\begin{itemize}
    \item After a GP or telehealth consultation, the prescription arrives as an eScript token
    \item The token is sent by SMS, email, or stored in a digital wallet
    \item The patient can choose a pharmacy, which is sent the prescription directly
    \item The pharmacist scans the token to dispense the medicine
    \item Repeats are available electronically
    \item A paper copy can be printed if the patient prefers
\end{itemize}

\section*{Differential Diagnosis}
eScripts are one of several ways medicines are prescribed.
\begin{itemize}
    \item Paper prescriptions, which are still used and accepted
    \item Electronic prescribing in hospital, where the hospital's system manages the medicines
    \item Medicines that are not prescribed, such as over-the-counter purchases
    \item The choice depends on the practice, the pharmacy, and patient preference
\end{itemize}

\section*{When to Seek Medical Care}
People should contact their doctor or pharmacist if an eScript has not arrived, if the token does not work, or if the wrong medicine or dose appears. Pharmacists can usually resolve problems with the token quickly.

\textbf{Contact your local emergency services for:}
\begin{itemize}
    \item A medicine emergency, such as running out of a critical medicine or a suspected overdose
    \item Severe symptoms requiring urgent medical care
    \item An allergic reaction to a medicine
\end{itemize}

\section*{Diagnosis and Investigations}
Using an eScript involves simple steps.
\begin{itemize}
    \item \textbf{Receiving:} the token arrives by SMS, email, or through the practice
    \item \textbf{Choosing a pharmacy:} the pharmacy may be selected online or the token taken to any participating pharmacy
    \item \textbf{Dispensing:} the pharmacist scans the code
    \item \textbf{Verification:} check the medicine and dose at collection
    \item \textbf{Repeats:} note how many repeats remain and their expiry
\end{itemize}

\section*{Management}
eScripts are managed through everyday pharmacy visits.
\begin{itemize}
    \item \textbf{Keeping the token:} store it in a safe place, such as a digital wallet
    \item \textbf{Using repeats:} the pharmacist records each repeat
    \item \textbf{Problems:} a lost or expired eScript can be resent by the doctor
    \item \textbf{Medication reviews:} bring your medicines, including eScripts, to reviews
    \item \textbf{Safety:} never share your token with others
\end{itemize}

\section*{Complications}
\begin{itemize}
    \item Lost or expired eScript tokens, requiring a new prescription
    \item Technical problems at the pharmacy
    \item Confusion about how many repeats remain
    \item Security concerns if the phone or token is shared
    \item People without smartphones or digital access may need paper scripts or help
\end{itemize}

\section*{Prognosis and Outcomes}
eScripts are convenient, secure, and reliable, and most people find them easier than paper prescriptions. They reduce lost scripts and pharmacy errors, and they work well with telehealth. The system continues to improve, and digital prescriptions are now a standard part of local health care.

\section*{Prevention and Self-Care}
\begin{itemize}
    \item Check that your contact details are current with your doctor so tokens arrive
    \item Store tokens in a secure digital wallet
    \item Check the medicine and dose when collecting
    \item Ask the pharmacist to explain any repeats
    \item Report any problem with an eScript promptly
    \item Ask for a paper script if you prefer or need one
\end{itemize}

\section*{Sources and Further Reading}
The following reputable sources provide further information for healthcare students and patients seeking reliable, current advice about electronic prescriptions.
\begin{itemize}
    \item Australian Digital Health Agency. \textit{Electronic prescriptions information.} https://www.digitalhealth.gov.au/
    \item Services Australia. \textit{Medicare and electronic prescriptions.} https://www.servicesaustralia.gov.au/
    \item Pharmacy Guild of Australia. \textit{Electronic prescriptions for patients.} https://www.guild.org.au/
    \item Healthdirect Australia. \textit{Getting a prescription.} https://www.healthdirect.gov.au/
    \item Department of Health and Aged Care. \textit{Electronic prescribing.} https://www.health.gov.au/
    \item NPS MedicineWise. \textit{Medicines information and safety.} https://www.nps.org.au/
\end{itemize}
